\section{Related Work}
\label{sec:related}

\paragraph{\textcolor{blue}{Secure update standards}}
\textcolor{blue}{The IETF SUIT working group specifies a firmware update architecture and a manifest information model for IoT devices~\cite{rfc9019,rfc9124}: a signed manifest binds the image digest, the device class it applies to, and a sequence number against rollback, so the image is protected regardless of how it is transported, and encrypting the image is optional.
TUF~\cite{samuel2010tuf} and Uptane~\cite{kuppusamy2016uptane} add role separation and signed, fresh metadata to survive key compromise and resist freeze and rollback attacks.
Baseline requirements ask for authenticated updates delivered over secure channels, and RFC~9325~\cite{rfc9325} gives current TLS recommendations.
Because these standards move integrity into the artefact, we assess each channel together with the image it carries.}

\paragraph{IoT traffic measurement and datasets}
{\emergencystretch=1em % avoids an overfull line at the unbreakable "Mon(IoT)r"
Ren et al.~\cite{10.1145/3355369.3355577} conducted a large-scale measurement of consumer IoT devices using the Mon(IoT)r testbed, covering 34,586 experiments across 81 devices; \textcolor{blue}{they classified encryption per flow, identifying TLS and QUIC by protocol, treating recognisable encoded or compressed content separately, and applying entropy thresholds only to the remaining flows. Their focus is information exposure across all traffic; they neither isolate update traffic nor} evaluate the security of the update channel itself\textcolor{blue}{, which is what we do with their dataset and our own captures}.
Dong et al.~\cite{10.1145/3618257.3624815} studied TLS usage and server-side certificate management across 2,014 real-world IoT devices sourced from a crowdsourced dataset, uncovering widespread use of customised TLS libraries and inconsistent PKI practices among vendors; however, their analysis is not restricted to update traffic and does not examine cipher-suite strength in the context of firmware delivery.\par}

\paragraph{IoT software update practices}
Prakash et al.~\cite{10.1145/3560835.3564551} inferred update practices from User-Agent strings in HTTP traffic, providing insight into versioning patterns but limited to cleartext channels.
Bradley and Barrera~\cite{10.1007/978-3-031-30122-3_25} characterised IoT update behaviour in a controlled lab environment, focusing exclusively on insecure HTTP exchanges without assessing encrypted channels or cryptographic configurations.
Dejon et al.~\cite{10.1007/978-3-030-41702-4_14} proposed an automated framework for the security analysis of IoT software updates, targeting binary-level properties of firmware images rather than the confidentiality or integrity of the update delivery channel.
Seitz et al.~\cite{icissp23} outlined secure software update requirements derived from industry standards, identifying gaps between specification and practice, though without empirical traffic measurement.

\paragraph{Firmware update security and vulnerabilities}
Wu et al.~\cite{294541} performed a systematic study of firmware update vulnerabilities across embedded devices, identifying flaws in update authenticity mechanisms; \textcolor{blue}{they analyse the device-side update logic in firmware binaries, whereas we observe the delivery channel and relate it to the protection of the delivered image, a complementary view of the same problem.}
Usman and Asplund analysed update-related vulnerabilities and proposed recommendations for reducing their exposure~\cite{10.1007/978-3-031-92886-4_7}, and characterised security risks in the update mechanisms of computing systems~\cite{USMAN2026105019}.
